\subsection{\texorpdfstring{分片 $W_p^1$ 函数的 Poincaré--Friedrichs 不等式}{分片 Wp1 函数的 Poincaré--Friedrichs 不等式}}
\label{sec:ch10-piecewise-poincare-friedrichs}
\setcounter{thmcount}{0}
\setcounter{equation}{0}

设 $\Omega\subset\mathbb R^2$ 是有界连通多边形区域, $1\leq p\leq\infty$. 本节目标是把下列 Poincaré--Friedrichs 不等式(参见习题~\MBUnresolvedReference{exr:ch04-source-4-x-4}{4.x.4} 和~\MBUnresolvedReference{exr:ch05-source-5-x-13}{5.x.13})推广到分片 $W_p^1$ 函数:
\MBSourceItem{1}
\begin{equation}
\MathBookFormulaID{formula-ch10-poincare-domain-mean}
\label{eq:ch10-poincare-domain-mean}
\|u\|_{L^p(\Omega)}\leq C\left(\left|\int_\Omega u\,dx\right|+|u|_{W_p^1(\Omega)}\right)
\qquad\forall u\in W_p^1(\Omega),
\end{equation}
\MBSourceItem{2}
\begin{equation}
\MathBookFormulaID{formula-ch10-poincare-boundary-mean}
\label{eq:ch10-poincare-boundary-mean}
\|u\|_{L^p(\Omega)}\leq C\left(\left|\int_{\partial\Omega}u\,ds\right|+|u|_{W_p^1(\Omega)}\right)
\qquad\forall u\in W_p^1(\Omega),
\end{equation}
其中正常数 $C$ 只依赖于 $\Omega$. 这类不等式可用于第~\ref{sec:ch10-nonconforming} 节和第~\ref{sec:ch10-discontinuous} 节中的有限元函数.

令 $\mathcal T$ 是 $\Omega$ 的单纯形三角剖分, $W_p^1(\Omega,\mathcal T)$ 表示关于 $\mathcal T$ 的分片 $W_p^1$ 函数空间, 即
\MBSourceItem{5}
\begin{equation}
\MathBookFormulaID{formula-ch10-piecewise-wp-space}
\label{eq:ch10-piecewise-wp-space}
W_p^1(\Omega,\mathcal T)
=\{v\in L^p(\Omega):v_T:=v|_T\in W_p^1(T)\quad\forall T\in\mathcal T\}.
\end{equation}
$\mathcal T$ 的内部边集和边界边集分别记为 $\mathcal E^i(\mathcal T)$ 和 $\mathcal E^b(\mathcal T)$, 顶点集记为 $\mathcal V(\mathcal T)$. 为把式~\eqref{eq:ch10-poincare-domain-mean} 和~\eqref{eq:ch10-poincare-boundary-mean} 推广到 $W_p^1(\Omega,\mathcal T)$, 先对分片常数函数空间 $\mathcal P_0(\Omega,\mathcal T)$ 建立 Poincaré--Friedrichs 不等式.

\MBSourceItem{6}
\begin{Lemma}
\label{lem:ch10-piecewise-constant-poincare}
设 $1\leq p<\infty$. 存在只依赖 $\mathcal T$ 最小角的正常数 $C$, 使得对所有 $c\in\mathcal P_0(\Omega,\mathcal T)$,
\MBSourceItem{7}
\begin{equation}
\MathBookFormulaID{formula-ch10-piecewise-constant-domain-mean}
\label{eq:ch10-piecewise-constant-domain-mean}
\|c\|_{L^p(\Omega)}\leq C\left[
\left|\int_\Omega c\,dx\right|
+\left(\sum_{e\in\mathcal E^i(\mathcal T)}|e|^{1-p}
\|[\![c]\!]\|_{L^p(e)}^p\right)^{1/p}\right],
\end{equation}
并且
\MBSourceItem{8}
\begin{equation}
\MathBookFormulaID{formula-ch10-piecewise-constant-boundary-mean}
\label{eq:ch10-piecewise-constant-boundary-mean}
\|c\|_{L^p(\Omega)}\leq C\left[
\left|\int_{\partial\Omega}c\,ds\right|
+\left(\sum_{e\in\mathcal E^i(\mathcal T)}|e|^{1-p}
\|[\![c]\!]\|_{L^p(e)}^p\right)^{1/p}\right],
\end{equation}
其中跳跃 $[\![c]\!]$ 如第~\ref{sec:ch10-nonconforming} 节所定义.
\end{Lemma}

\MathBookSourcePage{307}
\begin{Proof}
给定 $c\in\mathcal P_0(\Omega,\mathcal T)$, 定义 $Ec$ 为关于 $\mathcal T$ 的连续分片线性函数, 它在每个顶点取 $c$ 的平均值:
\[
\MathBookFormulaID{formula-ch10-averaging-operator-ec}
(Ec)(x)=\frac1{|\mathcal T_x|}\sum_{T\in\mathcal T_x}c_T
\qquad\forall x\in\mathcal V(\mathcal T),
\]
其中 $\mathcal T_x$ 是 $\mathcal T$ 中共享顶点 $x$ 的三角形集合. 任取 $T\in\mathcal T$, 记 $\mathcal V_T$ 为其顶点集. 由 $Ec$ 的定义、Hölder 不等式和网格最小角条件,
\[
\MathBookFormulaID{formula-ch10-ec-local-difference}
\|c-Ec\|_{L^p(T)}^p
\leq Ch_T^2\sum_{x\in\mathcal V_T}\frac1{|\mathcal T_x|}
\sum_{T'\in\mathcal T_x}|c_T-c_{T'}|^p.
\]
任意 $T'\in\mathcal T_x$ 可通过 $\mathcal T_x$ 中相邻三角形的链连接到 $T$. 由 Hölder 不等式和跳跃的定义,
\[
\MathBookFormulaID{formula-ch10-triangle-chain-jump}
|c_T-c_{T'}|^p\leq|\mathcal T_x|^{p-1}
\sum_{e_j}|e_j|^{-1}\|[\![c]\!]\|_{L^p(e_j)}^p.
\]
合并以上估计得到
\MBSourceItem{9}
\begin{equation}
\MathBookFormulaID{formula-ch10-ec-lp-difference}
\label{eq:ch10-ec-lp-difference}
\|c-Ec\|_{L^p(\Omega)}^p
\leq C\left(\max_{x\in\mathcal V(\mathcal T)}|\mathcal T_x|\right)^{p-2}
\sum_{e\in\mathcal E^i(\mathcal T)}|e|\|[\![c]\!]\|_{L^p(e)}^p.
\end{equation}
由式~\eqref{eq:ch10-ec-lp-difference} 和引理~\MBUnresolvedReference{lem:ch04-source-4-5-3}{4.5.3}, 还有
\MathBookSourcePage{308}
\MBSourceItem{10}
\begin{equation}
\MathBookFormulaID{formula-ch10-ec-wp-difference}
\label{eq:ch10-ec-wp-difference}
\left(\sum_{T\in\mathcal T}|c-Ec|_{W_p^1(T)}^p\right)^{1/p}
\leq C\left(\max_{x\in\mathcal V(\mathcal T)}|\mathcal T_x|\right)^{1-2/p}
\left(\sum_{e\in\mathcal E^i(\mathcal T)}|e|^{1-p}
\|[\![c]\!]\|_{L^p(e)}^p\right)^{1/p}.
\end{equation}
现在由式~\eqref{eq:ch10-poincare-domain-mean}、三角不等式、Hölder 不等式以及式~\eqref{eq:ch10-ec-lp-difference}--\eqref{eq:ch10-ec-wp-difference},
\begin{align*}
\MathBookFormulaID{formula-ch10-piecewise-constant-domain-proof}
\|c\|_{L^p(\Omega)}
&\leq\|Ec\|_{L^p(\Omega)}+\|c-Ec\|_{L^p(\Omega)}\\
&\leq C\left[\left|\int_\Omega c\,dx\right|
+\left(\sum_{e\in\mathcal E^i(\mathcal T)}|e|^{1-p}
\|[\![c]\!]\|_{L^p(e)}^p\right)^{1/p}\right].
\end{align*}
这里还用了 $\max_x|\mathcal T_x|^{1-2/p}$ 可由只依赖最小角的常数控制. 这证明了式~\eqref{eq:ch10-piecewise-constant-domain-mean}.

对式~\eqref{eq:ch10-poincare-boundary-mean} 作同样论证. 边界差异项满足
\MathBookSourcePage{309}
\[
\MathBookFormulaID{formula-ch10-ec-boundary-difference}
\left|\int_{\partial\Omega}(Ec-c)\,ds\right|^p
\leq C\left(\max_x|\mathcal T_x|\right)^{p-2}
\sum_{e\in\mathcal E^i(\mathcal T)}\|[\![c]\!]\|_{L^p(e)}^p.
\]
结合这一估计和式~\eqref{eq:ch10-ec-wp-difference}, 即得式~\eqref{eq:ch10-piecewise-constant-boundary-mean}.
\end{Proof}

由引理~\ref{lem:ch10-piecewise-constant-poincare} 可把 Poincaré--Friedrichs 不等式推广到 $W_p^1(\Omega,\mathcal T)$. 定义平均算子
\[
\MathBookFormulaID{formula-ch10-edge-average-operator}
\mathcal M_ev=\frac1{|e|}\int_ev\,ds
\qquad\forall v\in L^1(e).
\]
Hölder 不等式给出
\MBSourceItem{11}
\begin{equation}
\MathBookFormulaID{formula-ch10-edge-average-bound}
\label{eq:ch10-edge-average-bound}
\|\mathcal M_ev\|_{L^p(e)}\leq\|v\|_{L^p(e)}
\qquad\forall v\in L^p(e),\quad1\leq p\leq\infty.
\end{equation}

\MBSourceItem{12}
\begin{Theorem}
\label{thm:ch10-piecewise-poincare}
设 $1\leq p<\infty$. 存在只依赖 $\mathcal T$ 最小角的正常数 $C$, 使得对所有 $v\in W_p^1(\Omega,\mathcal T)$,
\MBSourceItem{13}
\begin{equation}
\MathBookFormulaID{formula-ch10-piecewise-poincare-domain}
\label{eq:ch10-piecewise-poincare-domain}
\begin{split}
\|v\|_{L^p(\Omega)}\leq C\bigg[&\left|\int_\Omega v\,dx\right|
+\left(\sum_{T\in\mathcal T}|v|_{W_p^1(T)}^p\right)^{1/p}\\
&+\left(\sum_{e\in\mathcal E^i(\mathcal T)}|e|^{1-p}
\|\mathcal M_e[\![v]\!]\|_{L^p(e)}^p\right)^{1/p}\bigg],
\end{split}
\end{equation}
以及
\MBSourceItem{14}
\begin{equation}
\MathBookFormulaID{formula-ch10-piecewise-poincare-boundary}
\label{eq:ch10-piecewise-poincare-boundary}
\begin{split}
\|v\|_{L^p(\Omega)}\leq C\bigg[&\left|\int_{\partial\Omega}v\,ds\right|
+\left(\sum_{T\in\mathcal T}|v|_{W_p^1(T)}^p\right)^{1/p}\\
&+\left(\sum_{e\in\mathcal E^i(\mathcal T)}|e|^{1-p}
\|\mathcal M_e[\![v]\!]\|_{L^p(e)}^p\right)^{1/p}\bigg].
\end{split}
\end{equation}
\end{Theorem}

\begin{Proof}
任取 $v\in W_p^1(\Omega,\mathcal T)$, 令 $\overline v\in\mathcal P_0(\Omega,\mathcal T)$ 在每个 $T$ 上等于 $v$ 的均值. 由引理~\ref{lem:ch10-piecewise-constant-poincare}、三角不等式和带尺度变换的 Poincaré 不等式,
\MathBookSourcePage{310}
\begin{align*}
\MathBookFormulaID{formula-ch10-piecewise-poincare-domain-proof}
\|v\|_{L^p(\Omega)}
&\leq\|v-\overline v\|_{L^p(\Omega)}+\|\overline v\|_{L^p(\Omega)}\\
&\leq C\bigg[\left|\int_\Omega v\,dx\right|
+\left(\sum_T|v|_{W_p^1(T)}^p\right)^{1/p}
+\left(\sum_{e\in\mathcal E^i(\mathcal T)}|e|^{1-p}
\|\mathcal M_e[\![v-\overline v]\!]\|_{L^p(e)}^p\right)^{1/p}\\
&\hspace{5em}+\left(\sum_{e\in\mathcal E^i(\mathcal T)}|e|^{1-p}
\|\mathcal M_e[\![v]\!]\|_{L^p(e)}^p\right)^{1/p}\bigg].
\end{align*}
由式~\eqref{eq:ch10-edge-average-bound}、带尺度变换的迹定理和 Poincaré 不等式,
\[
\MathBookFormulaID{formula-ch10-piecewise-average-jump-bound}
\left(\sum_{e\in\mathcal E^i(\mathcal T)}|e|^{1-p}
\|\mathcal M_e[\![v-\overline v]\!]\|_{L^p(e)}^p\right)^{1/p}
\leq C\left(\sum_{T\in\mathcal T}|v|_{W_p^1(T)}^p\right)^{1/p}.
\]
于是式~\eqref{eq:ch10-piecewise-poincare-domain} 成立. 类似地, 使用式~\eqref{eq:ch10-piecewise-constant-boundary-mean}, 并由迹定理控制 $\int_{\partial\Omega}(v-\overline v)\,ds$, 得式~\eqref{eq:ch10-piecewise-poincare-boundary}.
\end{Proof}

\MathBookSourcePage{311}
令定理~\ref{thm:ch10-piecewise-poincare} 中 $p\uparrow\infty$, 得到:

\MBSourceItem{15}
\begin{Theorem}
\label{thm:ch10-piecewise-poincare-infinity}
存在只依赖 $\mathcal T$ 最小角的正常数 $C$, 使得对任意 $v\in W_\infty^1(\Omega,\mathcal T)$,
\[
\MathBookFormulaID{formula-ch10-piecewise-poincare-infinity-domain}
\|v\|_{L^\infty(\Omega)}\leq C\left[
\left|\int_\Omega v\,dx\right|+\max_{T\in\mathcal T}|v|_{W_\infty^1(T)}
+\max_{e\in\mathcal E^i(\mathcal T)}|e|^{-1}
\|\mathcal M_e[\![v]\!]\|_{L^\infty(e)}\right],
\]
\[
\MathBookFormulaID{formula-ch10-piecewise-poincare-infinity-boundary}
\|v\|_{L^\infty(\Omega)}\leq C\left[
\left|\int_{\partial\Omega}v\,ds\right|+\max_{T\in\mathcal T}|v|_{W_\infty^1(T)}
+\max_{e\in\mathcal E^i(\mathcal T)}|e|^{-1}
\|\mathcal M_e[\![v]\!]\|_{L^\infty(e)}\right].
\]
\end{Theorem}

作为应用, 考虑关于 $\mathcal T$ 的非协调 $\mathcal P_1$ 有限元函数 $v$, 并假设 $\int_\Omega v\,dx=0$, 或 $v$ 在边界边的中点为零. 直接应用定理~\ref{thm:ch10-piecewise-poincare}、\ref{thm:ch10-piecewise-poincare-infinity} 和中点公式, 得
\[
\MathBookFormulaID{formula-ch10-nonconforming-poincare-finite-p}
\|v\|_{L^p(\Omega)}\leq C\left(\sum_{T\in\mathcal T}|v|_{W_p^1(T)}^p\right)^{1/p}
\qquad1\leq p<\infty,
\]
以及
\[
\MathBookFormulaID{formula-ch10-nonconforming-poincare-infinity}
\|v\|_{L^\infty(\Omega)}\leq C\max_{T\in\mathcal T}|v|_{W_\infty^1(T)},
\]
其中 $C$ 只依赖 $\mathcal T$ 的最小角.

最后把 Poincaré--Friedrichs 不等式推广到关于 $\Omega$ 的单纯形分割 $\mathcal P$ 的分片 $W_p^1$ 函数, 从而可用于第~\ref{sec:ch10-discontinuous} 节的间断有限元函数. 定义
\[
\MathBookFormulaID{formula-ch10-piecewise-wp-partition-space}
W_p^1(\Omega,\mathcal P)=\{v\in L^p(\Omega):v_T:=v|_T\in W_p^1(T)
\quad\forall T\in\mathcal P\}.
\]
$\mathcal P$ 的边集和顶点集如第~\ref{sec:ch10-discontinuous} 节定义, 内部边集记为 $\mathcal E^i(\mathcal P)$. 用 $\mathcal P$ 的最小角以及
\[
\MathBookFormulaID{formula-ch10-partition-rho}
\varrho(\mathcal P)=\max\{|\partial T|/|e|:T\in\mathcal P,
e\in\mathcal E(\mathcal P),\ e\subset\partial T\}
\]
度量 $\mathcal P$ 的正则性.

\MathBookSourcePage{312}
\MBSourceItem{16}
\begin{Theorem}
\label{thm:ch10-piecewise-poincare-partitions}
存在只依赖 $\mathcal P$ 的最小角和 $\varrho(\mathcal P)$ 的正常数 $C$, 使得对任意 $v\in W_p^1(\Omega,\mathcal P)$ 和 $1\leq p<\infty$, 定理~\ref{thm:ch10-piecewise-poincare} 的不等式~\eqref{eq:ch10-piecewise-poincare-domain} 和~\eqref{eq:ch10-piecewise-poincare-boundary} 在把 $\mathcal T$ 替换为 $\mathcal P$ 后仍成立; 对任意 $v\in W_\infty^1(\Omega,\mathcal P)$, 定理~\ref{thm:ch10-piecewise-poincare-infinity} 的两个不等式在作同样替换后仍成立.
\end{Theorem}

\begin{Proof}
连接 $\mathcal P$ 中每个三角形的中心与属于其边界的 $\mathcal P$ 顶点, 得到 $\Omega$ 的单纯形三角剖分 $\mathcal T$. $\mathcal T$ 的最小角只依赖于 $\mathcal P$ 的最小角和 $\varrho(\mathcal P)$. 又因为 $W_p^1(\Omega,\mathcal P)\subset W_p^1(\Omega,\mathcal T)$, 且 $v\in W_p^1(\Omega,\mathcal P)$ 跨越不属于 $\mathcal P$ 的任何 $\mathcal T$ 边的跳跃为零, 所以结论直接由定理~\ref{thm:ch10-piecewise-poincare} 和~\ref{thm:ch10-piecewise-poincare-infinity} 得到.
\end{Proof}

\MBSourceItem{17}
\begin{Remark}
\label{rem:ch10-poincare-without-edge-average}
由式~\eqref{eq:ch10-edge-average-bound}, 定理~\ref{thm:ch10-piecewise-poincare}、\ref{thm:ch10-piecewise-poincare-infinity} 和~\ref{thm:ch10-piecewise-poincare-partitions} 中的不等式去掉算子 $\mathcal M_e$ 后仍然成立.
\end{Remark}

\MBSourceItem{18}
\begin{Remark}
\label{rem:ch10-poincare-general-partitions}
本节的 Poincaré--Friedrichs 不等式可以推广到一般分割和三维区域(见 Brenner 2003a). 它们在间断 Galerkin 方法分析中的应用, 例如可见 Cockburn, Kanschat and Schötzau (2005), Carrero, Cockburn and Schötzau (2006), Gudi and Pani (2007), 以及 Brenner and Owens (2007).
\end{Remark}
